A system $\mathfrak{a}$ of integers in a number field $K$ is called an ideal if:
\begin{enumerate}
\item[(1)] Along with $\alpha$ and $\beta$, $\mathfrak{a}$ always contains $\alpha + \beta$ and $\alpha - \beta$ (it is a module).
\item[(2)] Along with $\alpha$ in $\mathfrak{a}$ and any integer $\lambda$ of $K$, the product $\lambda\alpha$ also belongs to $\mathfrak{a}$.
\end{enumerate}
The ideal consisting of the single number $0$ is denoted by $(0)$. An ideal of the form $(\alpha_1, \ldots, \alpha_r)$, the set of all $\sum \lambda_i \alpha_i$ with integral $\lambda_i$, is said to be generated by $\alpha_1, \ldots, \alpha_r$.
